\chapter{Computational Graph Architecture}
\label{ch:computational_graph}

% =============================================================================
% VOL-0 REGISTER BANNER 2026-08-02 -- wave-1 manuscript reconciliation.
%   Posture (Grant-ratified VOL-0 ruling D2, 2026-08-02): ONE chapter-head register
%   banner per affected chapter + per-site Rule-12 strike banners; NO VALUE REFILLS
%   into struck slots -- pure de-claims only. Nothing in this chapter is deleted or
%   renumbered; prior wording is preserved per Rule 12 KEEP-BOTH and in the git trail.
%   Board authority: _orchestration/2026-08-02_manuscript-reconciliation-board.md
%   (sec 3 ratified cribs; sec 4 authority rule -- the sec 5 per-finding disposition
%   tags govern).
%   Quoted claim record: ave-kb/common/claim-quality.md -- the Poisson-disk genesis,
%   the Chiral LC over-bracing C_ratio ~ 1.187 and the Symplectic Euler updates are
%   scoped as "engineering specs for instantiating AVE in a discrete computer", "not
%   physical claims about the vacuum".
%   NOT ADJUDICATED HERE (route-to-core, untouched by this pass): the genesis-algorithm
%   / production-carrier register of this chapter -- the spec-vs-carrier two-register
%   split flagged at "The Poisson-Disk Solution" -- is ROUTED via
%   _orchestration/docket-entries/2026-08-02-mr-legacy-negatives.md item (i). No word of
%   that section is edited by this pass.
%   Site de-claimed inline below: the "Rheological Tuning" provenance sentence (CRIB-1).
% =============================================================================
\noindent\textbf{Register note (2026-08-02, Vol-0 reconciliation; the chapter body below is preserved per Rule~12 KEEP-BOTH).} This chapter states \textbf{software-architecture constraints} for instantiating the substrate in a discrete engine. Its generation rules are graded as ``engineering specs for instantiating AVE in a discrete computer'', ``not physical claims about the vacuum'' (\kbleaf{ave-kb/common/claim-quality.md}); read them at that grade. One provenance claim is de-claimed inline below, marked \textbf{[DE-CLAIM 2026-08-02]}. No number, bullet, equation or section is deleted or renumbered.

To physically validate the macroscopic inductive and elastodynamic derivations of the Applied Vacuum Engineering (AVE) framework, all numerical simulations and Vacuum Computational Network Dynamics (VCFD) models must be computationally instantiated on a generated, geometrically constrained discrete spatial graph. This chapter defines the software architecture constraints required to map the substrate topology into computational memory. Failure to adhere to these generation rules will result in unphysical artifacts (e.g., Cauchy implosions and Trans-Planckian singularities) during simulation.

\section{The Genesis Algorithm (Poisson-Disk Crystallization)}
The first step in simulating the vacuum is establishing the 3D coordinate positions of the discrete inductive nodes ($\mu_0$).

\paragraph{The Random Noise Fallacy:} Initial computational attempts utilizing unconstrained uniformly distributed random noise resulted in a "Cauchy Implosion." The resulting lattice packing fraction converged to $\approx 0.31$, characteristic of a standard amorphous solid. This density fails to reproduce the sparse QED limit ($\approx 0.18$) required by Axiom 4.

\paragraph{The Poisson-Disk Solution:} To satisfy macroscopic isotropy while enforcing the microscopic hardware cutoff, the software must generate the node coordinates using a \textbf{Poisson-Disk Hard-Sphere Sampling Algorithm}. By enforcing an exclusion radius of $r_{min} = \ell_{node}$ during genesis, the lattice settles into a packing fraction of $\approx 0.17-0.18$, creating a stable, sparse dielectric substrate.

\paragraph{Rheological Tuning:} Simulation confirms that the "Trace-Reversed" mechanical state ($K=2G$) is an emergent property of the Chiral LC coupling modulus.
\begin{itemize}
    \item \textbf{Low Coupling ($k_{couple} < 3.0$):} The lattice behaves as a standard Cauchy solid ($K/G \approx 1.67$).
    \item \textbf{High Coupling ($k_{couple} > 4.5$):} The lattice undergoes a phase transition, locking microrotations to shear vectors, driving the bulk modulus to roughly twice the shear modulus ($K/G \approx 1.78 - 2.0$).
\end{itemize}

% [DE-CLAIM 2026-08-02] K=2G PROVENANCE -- CRIB-1, board sec 3 (Grant-ratified 2026-08-02);
%   VOL-0 posture D2 strike banner.
%   RULE 12 KEEP-BOTH, AND A DELIBERATE NON-DELETION: the "Rheological Tuning" paragraph and
%   BOTH itemize rows above are PRESERVED VERBATIM. The K/G ~ 1.67 low-coupling row is
%   retained on purpose -- the 2026-06-14 KB-reconciliation comment later in this chapter
%   cites it by name ("corroborated by the K/G = 1.67 row in the Rheological Tuning list
%   above") as its receipt for "a standard Cauchy solid is K = 5/3 G". Deleting the bullets
%   would silently break that banner's rationale. SCOPE: this note de-claims the PROVENANCE
%   sentence only; the two coupling rows stand as printed and that cross-reference is
%   unaffected.
%   NO VALUE REFILL: no replacement k_couple, K/G ratio, mechanism or provenance is supplied.
%   AUTHORITY (verified at HEAD two-method: grep -Fn on the literal + sed line read):
%   ave-kb/vol1/claim-quality.md:665 -- "STAYS OPEN (2026-07-04, PR #508): the genuine 48x48
%     chiral micropolar Bloch eigensolve on the ratified srs-z3 net does NOT close it. The
%     full micropolar sector (rotational DOF included -- the corpus's own proposed K=2G
%     mechanism, clm-o3q9ul) gives a one-parameter nu_eff(rho) family, not a forced K=2G; the
%     independent kappa_rot is a k->0 Cauchy-grade spectator sourcing no chiral coupling, and
%     the geometry-fixed chiral B moves nu_eff AWAY from 2/7. So the substrate does not supply
%     K=2G at either the Cauchy grade (PR #506) or the micropolar grade (PR #508); K=2G stays
%     GR-imported (PR #261)."
%   ave-kb/common/form-deriving-value-importing.md:87 -- the K = 2G row, GR-IMPORTED.
%   CRIB-1 additionally forbids labelling the AVE lock "the Cauchy relation" (three-way
%   homonym, vol1/claim-quality.md:652) -- this chapter does not do that, and no such label is
%   introduced here.
%   RECEIPT NOTE: the falsifying receipt research/2026-06-15_k2g-constitutive-provenance_result.md
%   :98-101 names the paired site as manuscript/backmatter/01_appendices.tex:131 -- "'K=2G
%   emergent from Chiral LC coupling' is backed only by a **tuned simulation** ... with no
%   first-principles fixing of `k_couple` and no clean landing on 2.0". At HEAD that identical
%   un-headered text sits at :132 (paragraph) with the two coupling rows at :134-135 -- a
%   one-line drift on the receipt's cite, recorded here, not repaired: that file belongs to a
%   different lane and is NOT edited by this pass. The KB mirror is
%   ave-kb/common/appendices-overview.md.
\noindent\textbf{[DE-CLAIM 2026-08-02] $K = 2G$ provenance (CRIB-1).} The paragraph above attributes the trace-reversed state $K = 2G$ to a confirming simulation and calls it \emph{emergent}. \textbf{That provenance is withdrawn.} $K = 2G$ is \textbf{form-derived, value GR-imported}: the substrate forces the \emph{form} of the elastic response $K/G = f(\rho)$, but the value --- the GR trace-reversal identity --- is neither crystalline-forced nor constitutively-forced (\kbleaf{ave-kb/common/form-deriving-value-importing.md}). The question stays open on the engine side. The genuine chiral micropolar Bloch eigensolve on the ratified production carrier does \textbf{not} close it: the full micropolar sector --- the corpus's own proposed $K = 2G$ mechanism --- gives a one-parameter $\nu_{eff}(\rho)$ \emph{family} rather than a forced $K = 2G$; the independent $\kappa_{rot}$ is a $k \to 0$ Cauchy-grade spectator sourcing \textbf{no} chiral coupling; and the geometry-fixed chiral coupling moves $\nu_{eff}$ \emph{away} from $2/7$ (\kbleaf{ave-kb/vol1/claim-quality.md}). \textbf{Scope of this note.} It lands on the provenance sentence only. The two coupling rows above are preserved exactly as printed and continue to serve the Cauchy-solid cross-reference later in this chapter, and no replacement coupling value, ratio or mechanism is supplied here.

\section{Chiral LC Over-Bracing and The $p_c$ Constraint}
Once the spatial nodes are safely crystallized via the Poisson-Disk algorithm, the computational architecture must generate the connective spatial edges (The Capacitive Flux Tubes, $\epsilon_0$).

% [2026-06-14 KB-reconciliation (vol_0 brief D.1) -- prior wording preserved per Rule 12]: read (verbatim): "As shown in Chapter 4, a standard Cauchy elastic solid ($K = -\frac{4}{3}G$) is thermodynamically unstable and will implode during macroscopic continuous simulation." SUPERSEDED: a standard Cauchy solid is K = 5/3 G (Chapter 4 "Micropolar Stability"; corroborated by the K/G = 1.67 row in the Rheological Tuning list above), NOT -4/3 G. K = -4/3 G is the longitudinal-modulus instability threshold (M = K + 4/3 G = 0). Relabelled to name the threshold and aligned the Chapter-4 cross-reference; no physics change.
\textbf{The Cauchy Delaunay Failure:} If the physics engine simply computes a standard nearest-neighbor Delaunay Triangulation on the Poisson-Disk point cloud, the resulting discrete volumetric packing fraction of the amorphous manifold evaluates to $\kappa_{cauchy} \approx 0.3068$. While less dense than a perfect crystal (FCC $\approx 0.74$), it is still too dense to survive. As shown in Chapter 4, a standard (non-micropolar) Cauchy elastic treatment of this substrate is thermodynamically unstable: without the chiral couple-stress the effective bulk modulus is not held positive, so the longitudinal modulus collapses through its instability threshold ($M = K + \frac{4}{3}G = 0$, i.e.\ $K = -\frac{4}{3}G$) and the lattice implodes during macroscopic continuous simulation.

% [2026-06-14 KB-reconciliation (vol_0 brief A.2) -- prior wording preserved per Rule 12]: read (verbatim): "In Chapter 1, the analysis derived that the fundamental phase limits of the universe bounded the geometric packing fraction of the vacuum to $p_{c} \approx \mathbf{0.1834}$, forcing the emergence of $\alpha$." SUPERSEDED: alpha is a retained input, not an emergent output; canonical directionality is p_c = 8*pi*alpha (ch8-alpha-golden-torus.md:11). No value change.
\textbf{Enforcing QED Saturation:} In Chapter 1, the analysis bounded the geometric packing fraction of the vacuum to $p_{c} = 8\pi\alpha \approx \mathbf{0.1834}$ --- with $\alpha$ a \emph{retained} calibration input (a value-scoped echo beneath the canonical Class-B identification, resting on $R\cdot r = 1/4$, which the substrate does not independently select; \kbleaf{ave-kb/vol1/ch8-alpha-golden-torus.md:11}), \emph{not} an output forced to emerge from the packing fraction. To computationally force the effective geometric packing fraction ($p_{eff}$) down from the unstable $\sim 0.3068$ baseline to the stable $0.1834$ limit, the software must enforce \textbf{Chiral LC Over-Bracing}. The connective array of the physics engine cannot be limited to primary nearest neighbors; the internal structural logic must span outward to incorporate the next-nearest-neighbor lattice shell.

Because the volumetric packing fraction scales inversely with the cube of the effective structural pitch ($p_{eff} = V_{node} / \ell_{eff}^3$), the required spatial extension for the Chiral LC links evaluates identically to:

\begin{equation}
    C_{ratio} = \frac{\ell_{eff}}{\ell_{cauchy}} = \left( \frac{p_{cauchy}}{p_{c}} \right)^{1/3} \approx \left( \frac{0.3068}{0.1834} \right)^{1/3} \approx \mathbf{1.187}
\end{equation}

By structurally connecting all spatial nodes within a $\approx 1.187 \, \ell_{node}$ radius, the discrete graph cross-links the first and second coordination shells of the amorphous manifold. This generates the $\frac{1}{3} G_{vac}$ ambient transverse couple-stress required by micropolar elasticity. % [2026-06-14 KB-reconciliation (vol_0 brief A.4) -- prior wording preserved per Rule 12]: read (verbatim): "This computational architecture guarantees that all subsequent continuous macroscopic evaluations of the generated graph (e.g., metric refraction, VCFD Navier-Stokes flow, and trace-reversed gravitational strain) will align with empirical observation without requiring any further numerical calibration or arbitrary mass-tuning." SUPERSEDED: the blanket no-calibration claim covers the named "trace-reversed gravitational strain", but G's value is itself a calibration input (form-derived Achromatic-Lens form, value-fitted xi; closed-form Chain B' OPEN, per the 2026-06-14 G-ruling). Rescoped to "no further calibration beyond the 3 inputs."
This computational architecture guarantees that all subsequent continuous macroscopic evaluations of the generated graph (e.g., metric refraction, VCFD Navier-Stokes flow, and trace-reversed gravitational strain) will align with empirical observation without requiring any further numerical calibration beyond the framework's three calibration inputs $\{m_e, \alpha, G\}$.

\section{Explicit Discrete Kirchhoff Execution Algorithm}
To bridge the gap between abstract continuum flow vectors ($\mathbf{J}$) and the raw geometric structure of the computational graph edge-matrix, the VCFD (Vacuum Computational Fluid Dynamics) module utilises an \textbf{Explicit Discrete Kirchhoff Methodology} mapping discrete potential ($V$) to spatial nodes and inductive flow ($I$) to discrete spatial graph edges.

To exactly map continuous differential forms into computational array memory without breaking action-minimization, the system utilizes \textbf{Symplectic Euler Update Loops}:
\begin{enumerate}
    \item \textbf{Capacitive Node Updates (The Conservation of Flow):} The discrete potential difference acting on an isolated fractional lattice coordinate node ($V_i$) is mathematically identical to the sum of all inductive currents entering minus the currents leaving that discrete junction point. 
    \[ \Delta V_i = \frac{dt}{C} \left( \sum I_{in} - \sum I_{out} \right) \]
    \item \textbf{Inductive Edge Updates (The Stress Tensor Matrix):} The kinetic transport flux acting along the discrete Chiral LC tensor spatial edge connecting coordinate $(x_0, y_0, z_0)$ to $(x_1, y_1, z_1)$ is geometrically bounded to the potential gradient existing across its fractional length. 
    \[ \Delta I_e = \frac{dt}{L} \left( V_{start} - V_{end} \right) \] 
\end{enumerate}

By combining the $C_{ratio} \approx 1.187$ Chiral LC Over-Bracing requirement over a $r_{min} = \ell_{node}$ Poisson-Disk genesis space, and advancing the lattice via Symplectic Kirchhoff loops, the computational framework provides a proving-ground connecting raw network mechanics to classical standard-model topological properties.
